\SPBGSplitSlide
  {Границы разделяющей полосы}
  {.330}{.50}
  {%
    \ReportSvmMargin
  }
  {.785}{.33}
  {%
    \SplitBody
    В нормированной записи рассматриваются три параллельные гиперплоскости:
    \[
      \begin{aligned}
        L_1&:\ \langle w,p\rangle+\beta=1,\\
        L_0&:\ \langle w,p\rangle+\beta=0,\\
        L_2&:\ \langle w,p\rangle+\beta=-1.
      \end{aligned}
    \]

    \begin{itemize}
  \item \(L_0\) -- разделяющая гиперплоскость;
  \item \(L_1\) и \(L_2\) -- границы разделяющей полосы.
\end{itemize}

    Можно доказать, что расстояние между границами равно
    \[
      h=\rho(L_1,L_2)=\frac{2}{\|w\|}.
    \]
  }
